\begin{apartats}

\apartat[1,25]{1}
Estudia la curvatura i troba els punts d'inflexió de la funció
\[
f(x)=x^4-6x^2+5 .
\]

\begin{solucio}
$f'(x)=4x^3-12x$ i $f''(x)=12x^2-12=12\left(x^2-1\right)$, que s'anul·la en $x=\pm1$.\\
$f''>0$ a $(-\infty,-1)$ i a $(1,+\infty)$: la funció és \textbf{convexa} (còncava cap amunt).\\
$f''<0$ a $(-1,1)$: hi és \textbf{còncava} (cap avall).\\
La curvatura canvia en tots dos punts, i $f(1)=f(-1)=0$: els punts d'inflexió són
$\boxed{(-1,0)}$ i $\boxed{(1,0)}$.
\end{solucio}

\begin{nomesllarg}
\apartat{0,75}
Estudia la curvatura de la funció $g(x)=x\,e^{x}$ i troba'n els punts d'inflexió.

\begin{solucio}
$g'(x)=e^{x}+x e^{x}=(x+1)e^{x}$ i $g''(x)=e^{x}+(x+1)e^{x}=(x+2)e^{x}$.\\
Com que $e^{x}>0$ sempre, el signe és el de $x+2$: \textbf{còncava} a $(-\infty,-2)$ i
\textbf{convexa} a $(-2,+\infty)$.\\
Punt d'inflexió: $\left(-2,-2e^{-2}\right)$.
\end{solucio}
\end{nomesllarg}

\apartat[1,25]{0,75}
Quants punts d'inflexió pot tenir, com a màxim, una funció polinòmica de grau $4$? Raona-ho.

\begin{solucio}
Els punts d'inflexió surten dels punts on s'anul·la la derivada segona i hi canvia de signe.
Si $f$ té grau $4$, aleshores $f''$ té grau $2$ i, com a màxim, dues arrels reals.\\
Per tant, \textbf{com a màxim dos}, com passa a l'apartat a). La funció $f(x)=x^4$ no en té
cap: $f''(x)=12x^2$ s'anul·la en $x=0$, però no hi canvia de signe.
\end{solucio}

\end{apartats}
